\section{Conclusion}
\label{sec:conclusion}

This study proposed endpoint-preserving compressed residual correction for a fixed base forecaster. The method combines a truncated DCT representation of the preceding residual block with its explicitly retained endpoint and coefficients of the current base forecast. Online ridge regressions estimate the correction within the original compressed output subspace. The contribution is the retained input construction, which preserves a recent residual value that the truncated coefficients do not generally determine.

The experiments support both the endpoint choice and the effectiveness of the complete method under the evaluated protocol. With the same regression dimensions and storage, the endpoint outperformed the first, middle, and mean residual summaries on all 12 seed-averaged dataset--backbone pairs under both warmup conditions. On the separate 72-condition legacy/refit evaluation, the complete method reduced MSE by an average of 16.99\% relative to the uncorrected base, with a median retained state of 4528 bytes. It also improved on the evaluated learned adapters when applied to the same joint base forecasts.

The complete input construction also reduced mean MSE by 4.71\% relative to endpoint-only regression on the 72 main conditions, although the benefit of residual DCT inputs was not universal. Setting sensitivity and prefix-based candidate selection retained positive average improvements on the previously used series. These retrospective checks do not supply an untouched validation set, and $K=4$ is an accuracy--storage choice rather than the most accurate tested setting.

The storage comparison places these gains within a tradeoff: full ELF achieved lower average error with substantially more retained arrays, while small adapters required fewer bytes. The local failure analysis further showed that full-period improvement can coexist with an adverse interval. These findings support explicit endpoint retention as an effective choice for the tested compressed correction model. Controlling local excess error remains a distinct objective, and retained array counts alone do not establish device performance.
